%!TEX root = ../template.tex
%%%%%%%%%%%%%%%%%%%%%%%%%%%%%%%%%%%%%%%%%%%%%%%%%%%%%%%%%%%%%%%%%%%%
%% chapter2.tex
%% NOVA thesis document file
%% State Of The Art
%%%%%%%%%%%%%%%%%%%%%%%%%%%%%%%%%%%%%%%%%%%%%%%%%%%%%%%%%%%%%%%%%%%%

\typeout{NT FILE chapter2.tex}%

\chapter{State Of The Art}
\label{cha:soa}

This chapter presents the State of The Art for current research for basketball activity recognition and fatigue monitoring systems. It 
overviews the current state of wearables, mainly which sensors are usually used, some communication protocols and their constraints. It 
also presents some types of data and algorithms relevant for this research. Finally it shows some previous work around this topic. This chapter 
presents the state of the art in basketball activity recognition and fatigue monitoring systems. It provides an overview of current wearable 
technologies, focusing on prevalent sensor types, communication protocols, and their inherent constraints. Additionally, it examines data 
characteristics and algorithms relevant for this research, concluding with a comprehensive review of existing literature and related work 
in the field.

\section{Wearables}
\label{sec:wearables}
A wearable is any smart device with sensors that can be worn, embedded in clothing, accessories or in the body or even adhered to or tattooed on the skin. There are many different 
types of Wearables such as Smart Socks, Smart Insoles, Smartwatches, Smart Rings amongst others. These devices are able to measure many different types of data such as position, pressure, 
brightness, velocity, heart rate and many others~\cite{reviewWearableTech}.

\begin{figure}[h]
    \centering
    \includegraphics[width=0.7\linewidth]{differentWearables.jpg}
    \caption{Different wearables~\cite{wearablesAndIOT}.}
    \label{fig:difWearables}
\end{figure}

Nowadays there are many different wearables that are able to collect the same type of data. With this in mind, the choice for the right device doesn't only revolve around what type of information 
it can collect. Portability, processing power, accessibility, battery life are some other important characteristics.


\subsection{Applications}
\label{sub:aplications}

Wearables provide useful data for many applications. Their portability, affordability, easy connection to a mobile phone and non-prescription nature of these devices
 makes them an easy choice for several uses. Amongst others the main ones are Health, Activity Recognition and Sports and Tracking/Localization~\cite{wearablesAndIOT}. 



\subsubsection{Health}
\label{ssub:health}

In the health sector wearables have gained a lot of importance over the last few years. With smart devices, doctors and patients can get a continuous health monitoring which lead to better diagnostics and disease and injury prevention. 
The data provided by the devices is also used to design personalized insoles, shoes and other objects that aid health care and injury prevention~\cite{injuryWalking}.
Besides this, wearables are also being used by people to track their own health without any prescription. 

Construction companies have also been investing in wearables to prevent injuries in their staff when working with heavy loads. In 2020 there were 619 million people affected by lower 
back pain globally and the number is estimated to increase to 843 million by 2050~\cite{lowerBack}.


\subsubsection{Tracking and Localization}
\label{ssub:trackingAndLocalization}

Knowing the position of a person or an animal is important in many occasions. From tracking animals to understand their behaviour to simply finding where a lost elder is, this category has many useful applications.

These devices tend to have a great battery life efficiency since they have to last long periods of time without being charged. 

\subsubsection{Activity Recognition and Sports}
\label{ssub:activityRecognitionAndSports}

Even though this section has a lot of overlapping technology with the health segment, its use cases go beyond the health category. Besides the day-to-day monitoring these devices can also recognize and monitor your physical activities.

With information about timings, amount of force, position and others, coaches and players can prevent injuries and improve their skills.

\subsubsection{Others}
\label{ssub:others}

Wearables are also used for education, gaming, virtual and augmented reality, law enforcement besides other applications.
While health, sports, and tracking are the primary drivers of the wearable market, the technology has seamlessly permeated numerous other sectors, offering innovative solutions for entertainment, professional duties, 
and daily lifestyle enhancements.

Wearables have revolutionized the entertainment industry by enabling highly immersive digital experiences. Beyond visual immersion, haptic feedback suits and smart gloves are increasingly popular, allowing users to 
physically "feel" digital objects, environmental effects, and interactions within a game or simulation.

The aforementioned~\gls{VR} and~\gls{AR} wearables are heavily utilized beyond entertainment for educational purposes. In medical schools, students use augmented reality headsets to visualize 3D human anatomy and perform simulated 
surgeries. In industrial settings, smart glasses provide workers with hands-free, real-time schematics and remote guidance while performing complex maintenance on machinery, significantly reducing human error and 
training time.

In the defence and security sectors, wearables are critical tools for situational awareness and personnel safety. Body-worn cameras are now standard equipment for law enforcement, providing accountability and real-time 
video streaming to command centers. Additionally, advanced tactical vests equipped with embedded physiological sensors and~\gls{GPS} trackers allow commanders to monitor the stress levels, vital signs, and exact locations 
of soldiers or first responders in hazardous environments.

Finally, wearables have become deeply integrated into everyday convenience and fashion. Smart rings and~\gls{NFC}-enabled wristbands facilitate secure, contactless payments and access control (such as unlocking smart doors 
or vehicles).

\subsection{Wearables in Basketball}
\label{sub:wearablesInBasketball}

Basketball was created in 1891 by a Canadian physical education professor named James Naismith in Springfield, Massachusetts as an indoor game to avoid the rain. 

Sensors initially became part of the game through the ball manufacturing process in the 1950s to improve safety and standardization of its pressure. In the 1980s the players started 
using blood oxygen and heart rate monitors in their training to improve strength and conditioning. The development of image sensors led to the recording and live broadcasting of basketball 
games in the 1990s and the 2000s. Displacement and acceleration sensors starting being used in the 2010s to monitor the players' shooting form and other movements. 
Wearable technology is rapidly evolving and it is the next big step for sports statistics and performance improvement~\cite{perspective}.

\begin{figure}[h]
    \centering
    \includegraphics[width=0.8\linewidth]{timeline.jpg}
    \caption{Evolution of Sensor Technology in Basketball (adapted from~\cite{perspective}).}
    \label{fig:sensorTimeline}
\end{figure}

Dave Brailsford popularized the concept "Marginal gains" which consists in improving every single detail, no matter how small in order to get better. The biggest obstacle in 
this philosophy is quantifying and identifying what we can improve. The answer to this distress is Wearables~\cite{impactWearable}.

For this purpose the technology used consists mostly of micro-electromechanical systems such as Inertial Measurement Units (\gls{IMU}), networked sensor suites and exoskeletons. 
Risk assessment shares some of its data with performance enhancement devices working with 
pressure sensors and~\glspl{IMU} that host several sensors (e.g., accelerometers, gyroscopes, and magnetometers)~\cite{biomechanicalPerfomance}. 





\section{Movement Recognition Sensors}
\label{sec:movement_recognition_wearablesT}

There are many different technologies capable of capturing similar data through different mechanisms. % devInsole
For acquiring movement related data the usual methods are through an~\gls{IMU}, image capture or some sort of pressure sensors.

Another way to collect information about player movements is through the muscles, for example, with an~\gls{EMG} sensor.


\subsection{\gls{IMU} - Inertial Measurement Units}
\label{ssec:imu_STA}

An~\gls{IMU} is a sensor system design to capture motion and orientation usually consisting of an accelerometer, a gyroscope and sometimes a magnetometer. These components allow the measurement of acceleration, angular rate and 
magnetic field respectively. When combined, they provide a good estimation of position, orientation and movement patterns. However these approximations suffer from drift over time, meaning that small errors accumulate 
and lead towards increasingly inaccurate readings. Due to its low cost, compact size, high sample rate and versatility,~\glspl{IMU} are especially valuable for wearable devices for activity/movement recognition practices.

\subsection{\gls{EMG} - Electromyography}
\label{ssec:emg_STA}

\gls{EMG} is a technique for evaluating and measure electrical activity generated by muscles. With the generated data researchers and practitioners can assess the muscles contribution to different movements, evaluate neuromuscular 
coordination and identify abnormal patterns that may indicate fatigue, compensation or injury risk.
There are two main types of~\gls{EMG}:

\begin{itemize}
  \item Needle~\gls{EMG} - A technique commonly used in neurology that where fine-wire electrodes are inserted into the muscle tissue leading to very precise recordings but is very invasive.
  \item Surface~\gls{EMG} - A non-invasive technique where the electrodes are placed over the skin.
\end{itemize}


\subsection{Plantar Pressure Wearables}
\label{ssec:plantar_pressure_wearables_STA}

Plantar pressure wearables are innovative devices designed to monitor and analyze the distribution of pressure across the soles of the feet. By embedding sensors into insoles, smart 
socks, or even shoes, these wearables capture detailed data on foot loading patterns, which is critical for understanding gait, balance, and overall foot health. Plantar pressure data 
offers valuable insights for applications in sports performance, injury prevention, rehabilitation, and the management of conditions such as diabetes, where pressure distribution can 
indicate areas at risk for ulcers or other complications. They are also useful for better footwear designs or orthosis~\cite{coimbra, embedding, devInsole, deepLearn}.

For gait monitoring through pressure there are several different types of sensors used in wearables that can be separated in five different categories~\cite{smartSocks}:

\begin{enumerate}
  \item Capacitive Sensors - Consisting in two conductive plates separated by an insulating layer. When force is applied, they develop a voltage variation.
  \item Resistive Sensors - These are the most common sensors. Their resistance varies with the applied pressure.
  \item Optoelectronic Sensors - Composed by a transmitter and a receiver (a laser or a LED and a photodiode). When a force is applied, the light that reaches the photodiode varies proportionally.
  \item Piezoresistive and Piezoelectric Sensors - These sensors use the  piezoelectric effect. They convert the changes in applied pressure to electrical current.
\end{enumerate}

As of today there are four different kinds of devices that can measure plantar pressure~\cite{smartSocks}:

\begin{enumerate}
  \item Platform Systems - Known as the most reliable method for measuring plantar pressure, offering high accuracy and detailed spatial resolution. Their main limitation is a lack of portability, restricting their use to controlled environments like labs and hospitals.
  \item In-Shoe Systems: These systems provide greater mobility and flexibility, with optimized circuit design, power efficiency, and communication technology, and they're more affordable than platform systems. They record plantar pressure inside the shoe and allow for extensive data collection across multiple steps, enabling long-term gait analysis indoors and outdoors. However, they are less accurate than platform-based measurements.
  \item Smart Wireless Insoles: Unlike platform and in-shoe systems, smart wireless insoles are cordless, eliminating the need for wires to connect sensors or data acquisition units. While in-shoe systems aren't ideal for extended outdoor use, smart insoles can be used both indoors and outdoors. They often include Bluetooth or WiFi for data transmission and a power source. However, the added elastic layer they create inside the shoe can be several millimetres thick, potentially distorting the accuracy of foot pressure data. They are also relatively expensive and not suited for daily wear.
  \item Smart Socks: These fabric-based systems integrate sensors directly into the textile, which avoids the challenges of the previous systems. However, they require handmade or complex manufacturing techniques, which is a drawback.
\end{enumerate}

In basketball, smart insoles and smart socks are particularly well-suited wearables due to their portability, flexibility, and adaptability to both indoor and outdoor environments. 
Unlike larger, stationary measurement systems, these wearable devices are attached to the athlete, providing continuous and real-time data regardless of the setting.

Smart insoles and socks, when embedded with pressure sensors, can track a player's foot pressure distribution, stability, and impact forces with each step, jump, or landing. This information 
is crucial for basketball players, as it helps to analyze their movement patterns, improve balance, reduce injury risk, and enhance overall performance. The insole's thin profile and 
ability to fit inside regular shoes also ensure they don't interfere with the player's comfort or mobility.

Similarly, smart socks, which incorporate sensors directly into the fabric, offer flexibility and comfort while capturing dynamic foot pressure and movement data. Their lightweight 
design and close contact with the foot provide detailed insights into the nuances of foot pressure changes during the quick movements in basketball.

Both smart insoles and socks also benefit from wireless connectivity, allowing data to be transmitted to a smartphone or tablet, enabling real-time monitoring and analysis. This 
flexibility and data accessibility are particularly valuable in basketball training and performance analysis, as players and coaches can assess and adjust techniques in real time, 
whether in an indoor or outdoor court.
%Muscle Sensors

\subsection{Restraints}
\label{ssec:restraints}

Wearables have to be light and comfortable since they should be unnoticeable and not disturb the user's activities. This means that these gadgets have low processing capacities 
and their sensors monitor the data at lower frequencies then non wearable devices. Due to their size they often use miniature sensors that can be sensitive to placement and alignment. 
Inaccurate placement, external interference, or device movement can cause errors, reducing the reliability of collected data. Additionally, in applications where precise measurements 
are essential, the smaller sensors in wearables can't match the accuracy of larger, lab-based equipment.

Wearables are often exposed to a wide range of environmental conditions, including sweat, water, dust, and physical impacts. Ensuring these devices remain durable and function properly 
under such conditions is a challenge, as components like sensors and circuits need to be robust enough to withstand daily use and environmental exposure.

Due to their lack of memory and processing power, the data collected from the devices has to be sent to another device for storage or further processing. The main communication protocols 
with low energy consumption are~\gls{BLE} and WiFi for close data transmission. These methods can be impacted by limited range, interference, and connectivity issues, 
especially in crowded or signal-dense environments, which can disrupt real-time data streaming.

Since these devices are quite small, one of their main restraints is the battery life. With every action taken by the device, whether it's collecting data with a sensor, processing that data or sending it 
somewhere else, energy is always spent. Depending on their purpose, wearables may have to last up until one year without being charged~\cite{challengesWearables}. To overcome this challenge 
some devices rely on energy harvesting components to charge the device such as solar power, motion energy, etc~\cite{embedding}.

Minding the others problems already mentioned, wearables have low security so they are relatively easy targets due to their poor encryption and protection~\cite{wearablesAndIOT}.
Due to the lack of privacy of the athlete's data and since wearable technology can possibly obstruct and injury someone, wearables still are'nt accepted in many sports 
(in federated play), basketball included. Even tho they are't allowed in games, there aren't any rules for practices and some teams are taking advantage of these 
devices~\cite{NBPAWearables, hawksWear}.


\subsection{Communication Protocols} %dá para aumentar
\label{ssec:communication_protocols}

Wearable technology relies heavily on efficient communication protocols to ensure seamless data exchange between devices, applications, and networks. These protocols define the rules 
for data transmission, enabling wearables to monitor health, track activity, and connect with other devices in real time. Given the constraints of wearables, communication protocols 
play a critical role in balancing performance with power efficiency.

For this project two kinds of protocols will be discussed: Bluetooth Low Energy and WiFi.


\subsubsection{BLE - Bluetooth Low Energy}
\label{sssec:bleT}

BLE is a short-range communication technology known for its low power consumption. Classified as a Wireless Personal Area Network (WPAN), BLE is widely used in 
various wearable applications, from diabetes monitoring devices to smartwatches.

BLE's primary design objectives are to enable low-power communication, particularly for peripheral devices. It is optimized for affordability and simplicity, aligning well with many 
developer requirements and contributing to its swift adoption.

The main difference between BLE and the classic Bluetooth is their transmission rates. BLE transmits at around 2000 Kbps while Bluetooth Classic transmits at around 2-3 Mbps making its energy consumption 
around two times bigger than BLE. This is the reason BLE is the usual choice for IOT devices. Besides this difference BLE works the same way as Bluetooth Classic~\cite{masterInsole, bachelorInsole}.

\subsubsection{WiFI}
\label{sssec:wifiT}

WiFi is a widely used wireless Local Area Network (LAN) technology, primarily aimed at connecting mobile devices to the internet. Its broad adoption means it's available in most homes, businesses, and many public 
spaces. WiFi uses variations of the IEEE 802.11 protocol depending on the generation, with the latest generation, WiFi 6 (802.11ax), approved in September 2024. Operating on the 2.4 GHz or 5 GHz 
frequency bands, WiFi 6 also supports wider channels, optimized mainly for video streaming.

WiFi delivers high data speeds and easy setup for access points, making it a mature and well-proven technology. However, its range is somewhat limited indoors due to high signal 
absorption, though range improves in open, line-of-sight environments. Due to its relatively high power consumption, WiFi is not ideal for direct communication in IoT devices~\cite{masterInsole, bachelorInsole}.


\section{Movements}
\label{sec:movements}

Humanmovement forms the baseline for any activity recognition system. At its core, it relies on the repetitive and predictable patterns of 
standard walking and running. However, in the context of a high-intensity sport like basketball, it's necessary to understand its explosive, 
multi-directional, and highly complex movements.

\subsection{Gait Cycle}
\label{ssec:gait_cycleT}

The Gait Cycle is the sequence of movements that one leg takes in each step when walking or running. It is the whole moving process from the first point of contact between 
a foot and the ground until it reaches the ground again (see~\autoref{fig:gaitCycle})~\cite{masterInsole, freewalker}.

\vspace{.5cm}
\begin{figure}[h]
    \centering
    \includegraphics[width=0.8\linewidth]{gaitCycle.jpg}
    \caption{The gait cycle of a leg with with a cycle time of 0.92s. 40 ms interval between lines~\cite{gaitBook}.}
    \label{fig:gaitCycle}
\end{figure}

The usual method for measuring Gait Cycle is with a platform and a multi-camera setup since it is the most precise and reliable. However, this method is very expensive and can only be 
applicable indoors in a secure environment. An insole or a smart sock aren't as reliable but they compensate by being cheaper, more comfortable for the user and can be worn in their activities.

Although the Gait Cycle varies from person to person (minding age, body type and others) and even in each step each person takes, there are many tendencies for normal gait.
With this in mind these patterns can be used to predict and understand what forces are applied in a gait cycle~\cite{masterInsole}.

\begin{figure}[h]
  \centering
  \includegraphics[width=0.7\linewidth]{footDir.jpg}
  \caption{The three foot's directions~\cite{gaitBook}.}
  \label{fig:footDir}
\end{figure}

There are three foot's directions(see~\autoref{fig:footDir}): Vertical which is the direction perpendicular to the foot's palm. Fore-aft that follows the line between the toe to the heel. And the 
lateral direction which goes from one side of the foot to the other~\cite{masterInsole}.

There are two different stages for this motion in slower gaits and an extra stage in faster gaits~\cite{porto}:

\begin{itemize}
  \item Stance Phase - time period when the foot is on the ground. This phase represents about 60\% of the normal gait cycle.
  \item Swing Phase - when the foot is in the air. It normally represents the remaining 40\% of the gait cycle.
  \item Flight Phase - this phase only occurs in faster gaits such as running. It refers to the moments when neither foot is in contact with the ground.
\end{itemize}

\subsection{Basketball Movements}
\label{ssec:baskebtall_movements}

Basketball involves many other distinct movements besides gait. These motions are very complex making them difficult to accurately identify 
but they are essential for injury prevention and for the performance enhancement analysis. In this project data will be collected from three groups of movements with increasing complexity.

\begin{figure}[h]
  \centering
  \includegraphics[width=0.8\linewidth]{basketPos1.jpg}
  \caption{Some basketball stances.}
  \label{fig:basketPos}
\end{figure}

\subsubsection{Basic Movements}
\label{sssec:basic_movements}

\begin{itemize}
  \item Sprint(fast gait) - Basketball players frequently sprint down the court, which involves a pattern of quick, alternating footfalls with consistent gait cycles on each foot.
  \item Defence Slide/Stance - The defensive slide requires the player to move laterally while maintaining a low stance.
  \item Rebound Jump - Rebounding requires explosive jumping to reach the ball, sometimes with minimal preparatory movement, other times with a box out.
\end{itemize}

\subsubsection{Intermediate Movements}
\label{sssec:intermediate_movements}

\begin{itemize}
  \item Pivot - A pivot involves planting one foot firmly while the other foot moves to change direction.
  \item Crossover Dribble - In a crossover dribble, the player shifts their weight rapidly from one foot to the other, using pressure changes to decelerate in one direction and quickly accelerate in the opposite direction.
  \item Change of Direction in Defence - A player changes his direction while sliding.
  \item Boxout - The act of securing a positing usually near the basket to deny the rebound for the opponent.
  \item Screen - Similar to the boxout but it is done anywhere on the court to force a defensive player to change his path.
  \item Layup - A layup is normally a two step shot and it requires a step-by-step increase in pressure, usually peaking in the final step before jumping.
  \item Cutting - Cutting involves a sharp change in direction with a great speed shift. It is an off ball movement.
\end{itemize}

\subsubsection{Advanced Movements}
\label{sssec:advanced_movements}

\begin{itemize}
  \item Jumpshot - The act of shooting the ball. The jump shot is a fundamental technique in basketball, characterized by a burst of force as they push off the ground to jump.
  \item Eurostep - This movement is a specific type of layup where the player takes the first step in one direction and the second step to another direction.
  \item Step Back - The step back is a Jump Shot where the player takes a step back before shooting the ball.
  \item Failed Landing (Incomplete layup) - Any type of fall characterized by the sudden abnormal muscle activity and high acceleration and velocities.
\end{itemize}

%\section{Relevant Data for Analysis}


\section{Data Analysis}
\label{sec:data_for_analysis}

By leveraging advanced data analysis techniques, such as machine learning algorithms and statistical modelling, pressure insole data can 
reveal subtle patterns and correlations that inform training adjustments, injury risk, and performance improvements. There are many 
different algorithms with different objectives. In this section some of them will be explained:

\subsection{Algorithms}
\label{ssec:data_algorithm}

\subsubsection{K-nearest-neighbours classification}
\label{sssec:k-nearest_neighbours}

The K-nearest-neighbours is one of the simplest supervised machine learning algorithms. It's a classification algorithm which uses proximity to assigning each point to the class with 
the most similar elements within the K-nearest-neighbours~\cite{bachelorInsole, porto}. 

As an example observe~\autoref{fig:k-near}. If someone would have to assign a class (green triangles, blue squares or red circles) to the cross in the figure, they would 
probably say it belongs to the green triangles since they are the closest objects. But if instead someone had to assign a class for the star, it would be a much harder task to the 
a naked eye. In this case two factors make that decision: Which are the closest objects (depending on the weight k) and how many of each class are there in that group~\cite{bachelorInsole}. 
In other words imagine a small circle with its center on the object that we want to classify and if it isn't touching other know objects imagine it growing until it does.

\begin{figure}[h]
  \centering
  \includegraphics[width=.8\linewidth]{k-near.jpg}
  \caption{K-nearest-neighbours example - adapted from~\cite{bachelorInsole}}
  \label{fig:k-near}
\end{figure}

 
\subsubsection{Dynamic Time Warping}
\label{sssec:dynamic_time_warping}

The dynamic Time Warping (DTW) algorithm transforms unsynchronized time series or time series with different lengths allowing a dynamic comparison between them. This algorithm calculates 
the best possible alignment between two different series dependent on time providing a measurement of how well aligned (or not) they are~\cite{bachelorInsole}.

\subsubsection{Regression analysis}
\label{sssec:regression_analysis}

Regression analysis' purpose is to find a continuous outcome variable, usually y, through one or more predictor variables, normally x (if only one). The simplest of these models is the 
linear regression model that relates a linear relationship between x and y. This equation can be expressed as~\cite{bachelorInsole}:

\begin{equation}
  y = \beta + \alpha \cdot x + \epsilon
  \label{eq:regAnalEQ}
\end{equation}

In this equation y is the output variable, $\beta$ is the starting point (x=0), $\alpha$ is the regression coefficient associated with the predictor variable x and $\epsilon$ is the residual error.
An extended version of this equation with n predictor variables is:

\begin{equation}
  y = \beta + \alpha_1 \cdot x_1 + \alpha_2 \cdot x_2 + \alpha_n \cdot x_n + \epsilon
  \label{eq:regAnalEQGeneral}
\end{equation}


\subsubsection{Recursive Feature Elimination}
\label{sssec:recursive_feature_elimination}

Recursive Feature Elimination (RFE) selects features by progressively narrowing down the set of available features. Each feature is assigned a weight, and those with the lowest absolute 
weights are removed from the set. In each iteration, a single feature is removed, and this process continues on the reduced set until the target number of features is reached~\cite{porto}.


\subsubsection{Principal Component Analysis}
\label{sssec:principal_component_analysis}

Principal Component Analysis (PCA) reduces the dimensionality of data by applying a linear transformation based on singular value decomposition. This transformation projects the data 
into a lower-dimensional space, creating a new set of variables that are ordered by their importance~\cite{porto}.


\subsubsection{Linear Discriminant Analysis}
\label{sssec:linear_discriminant_analysis}

The Linear Discriminant Analysis (LDA) model fits a Gaussian distribution to each class, assuming a shared covariance matrix across all classes. This model reduces input dimensionality 
by projecting data onto the directions that best separate the classes. Its goal is to enhance separation between classes by maximizing the scatter between them while minimizing 
scatter within each class~\cite{porto, ComprehensiveStudyOnEMG}.


\subsubsection{Neural Network - NN}
\label{sssec:neural_network}

A~\gls{NN} is a~\gls{ML} model that imitates the structure and functions of biological neural networks. A~\gls{NN} utilizes many connected nodes called "artificial neurons" organized by 
layers where each node connects to others from the previous or the next layer (Figure~\autoref{fig:NNFigure}). Each of these connections have corresponding weights determined 
trough the training algorithm. For this process the data has to be labelled.


\begin{figure}[h]
  \centering
  \includegraphics[width=\linewidth]{neuralNetwork.jpg}
  \caption{Neural Network Example}
  \label{fig:NNFigure}
\end{figure}

There are three types of layers: input, hidden and output. A NN must have at least one layer of each, though it may contain multiple hidden layers. When a network has two or more hidden 
layers it is considered a Deep Neural Network (DNN). 
This~\gls{ML} technique is highly effective for analysing large datasets with non-linear relationships. However, to achieve accurate and precise results it requires a substantial 
amount of complete data, as missing entries can significantly affect performance. 
Despite its effectiveness, one major limitation of neural networks is their lack of interpretability, making it difficult to clearly understand the underlying classification 
boundaries~\cite{ComprehensiveStudyOnEMG}.

\subsubsection{Convolutional Neural Network - CNN}
\label{sssec:convolutional_neural_network}

A~\gls{CNN} is a sort of Deep NN designed to process spatial or temporal structures. These models, before the fully connected layers, use convolutional filters to 
automatically extract local features, reducing the number of parameters compared to fully connected networks improving efficiency.
CNNs perform well in large-scale classification tasks and require minimal manual feature engineering. However, they often need substantial computational resources and can be difficult 
to interpret. Despite these limitations, CNNs remain one of the most effective tools for pattern recognition in modern machine learning~\cite{Review_IRBM_EMG}.


\subsubsection{Fuzzy Logic - FL} 
\label{sssec:fuzzy_logic}

Fuzzy Logic is a form of multi-valued logic in which its values are real and range between 0 and 1. Providing a mathematical framework capable of incorporating and balancing both precision and significance in 
decision-making processes.

A FL system operates through two main processes: fuzzification and defuzzification.
Fuzzification converts crisp input values into fuzzy values by assigning them degrees of membership within a fuzzy set, typically represented by an arbitrary curve that defines 
the relationship between the input values and their corresponding membership levels. A fuzzy set $A$ can be expressed as:

\begin{equation}
  A = {(x,  \mu A(x))  |  x  \epsilon  U} 
  \label{eq:fuzzyLogicEQ}
\end{equation}

where $A$ is the fuzzy set, $U$ is the universe of discourse containing the elements $x$, and $\mu A(x)$ is the membership function.

After fuzzification, the fuzzy sets are evaluated through a series of if-then rules using logical operators to produce an output. This output is then converted back into crisp 
values through defuzzification, which can be performed using various methods such as the centroid or bisector approaches.
The main advantage of FL lies in its flexibility and its ability to be easily modified or combined with other classification techniques. However, FL also presents limitations - 
its numerous localized parameters and training requirements can make system construction and tuning a time-consuming process~\cite{ComprehensiveStudyOnEMG}.

\subsubsection{Neuro-Fuzzy Hybridization - NF} 
\label{sssec:neuro_fuzzy_hybridization}

NF combines the principles of FL and NN creating an intelligent hybrid model that merges human-like reasoning with adaptive learning. Typically, a NF system is built upon a fuzzy 
logic framework and trained using data-driven methods from NN theory. During training, it adjusts local parameters based on local data and represents knowledge as fuzzy rules, 
allowing initialization with or without prior information. This hybrid approach leverages the strengths of both processes, interpretability from FL and learning ability from NN, 
while reducing their main drawbacks such as dependence on prior knowledge and high computational demands~\cite{ComprehensiveStudyOnEMG}.


%\subsubsection{Support Vector Machine - SVM} % ComprehensiveStudyOnEMG
%\label{sssec:support_vector_machine}



\subsection{EMG Features}
\label{ssec:emg_features}
The raw~\gls{EMG} signal is very noisy and very high frequency so it generates large volumes of data making it very demanding for direct analysis making it inadequate for real-time 
processing or~\gls{ML} applications. Therefore, relevant features of the signal must be extracted to capture the essential characteristics of the muscle activity to reduce 
computational load, improve robustness against noise and improve the performance of classification algorithms. 

\subsubsection{Mean Absolute Value - MAV} 
\label{sssec:mean_absolute_value}

The Mean Absolute Value is one of the most common and simpler features extracted from~\gls{EMG} signals. It corresponds to the average of the absolute values of the signal over a 
specific period of time, providing a measurement of the muscle's electrical activity level. 
Mathematically, MAV is calculated as:

\begin{equation}
  MAV = \frac{1}{N} \sum_{i = 1}^{N} \left\lvert x_i \right\rvert 
  \label{eq:mavEQ}
\end{equation}

where $N$ is the number os samples in the time frame and $x_i$ is the~\gls{EMG} signal at sample $i$~\cite{HandGestureRecognition, ComprehensiveStudyOnEMG, 
BiomedicalSignalProcessingAndControl, Review_IRBM_EMG, ReviewOfEMGSignalClassification, Hybrid, ComparativeStudyOfSEMGFeatureEvaluationMethods, EvaluationOfFeatureExtractionAndClassification}.

%\subsubsection{Mean Absolute Value Slope - MAVS} % ComprehensiveStudyOnEMG
%\label{sssec:mean_absolute_value_slope}


\subsubsection{Root Mean Square - RMS} 
\label{sssec:root_mean_square}

The Root Mean Square is another common feature in EMG analysis. As the name implies, it is the square root of the mean of the signal amplitude emphasizing larger signal amplitudes 
than the MAV. It's defined as:

\begin{equation}
  RMS = \sqrt{\frac{1}{N} \sum_{i = 1}^{N} x_i^2}
  \label{eq:rmsEQ}
\end{equation}

Where $N$ is the number os samples in the time frame and $x_i$ represents the~\gls{EMG} signal at each sample~\cite{HandGestureRecognition, BiomedicalSignalProcessingAndControl, 
MeasuringSurfaceElectromyographyWithTextileElectrodes, Review_IRBM_EMG, ReviewOfEMGSignalClassification, Hybrid, EvaluationOfFeatureExtractionAndClassification, ComparativeStudyOfSEMGFeatureEvaluationMethods, 
ResearchOnLowerLimbMotionRecognition}.

\subsubsection{Variance - VAR} 
\label{sssec:variance}

The Variance of a signal measures how much it fluctuates around its mean, reflecting the degree of variability or dispersion in muscle activation. 

Mathematically it is expressed as:

\begin{equation}
  VAR = \frac{1}{N - 1} \sum_{i = 1}^{N} (x_i - \mu)^2
  \label{eq:varEQ}
\end{equation}

Where $N$ is the number of samples, $x_i$ represents each signal value, and $\mu$ is the mean of the signal.

It is a useful feature for pattern recognition and classification since it captures the power and intensity fluctuations of the signal helping differentiating the many types of 
muscle activity~\cite{HandGestureRecognition, ComprehensiveStudyOnEMG, Review_IRBM_EMG, ComparativeStudyOfSEMGFeatureEvaluationMethods, EvaluationOfFeatureExtractionAndRecognition, 
EvaluationOfFeatureExtractionAndClassification}.

\subsubsection{Waveform Length - WL}
\label{sssec:waveform_length}

The Waveform Length represents the cumulative length of the signal waveform over a given time frame. It reflects both the signal's complexity and amplitude variation, capturing 
information about the frequency and intensity of the muscle activity.

The WL feature is calculated according to:

\begin{equation}
  WL = \sum_{i = 1}^{N - 1} \lvert x_{i + 1} - x_i \rvert
  \label{eq:wlEQ}
\end{equation}

Where $x_i$ represents the $i$-th sample of the EMG signal, and $N$ is the total number of samples.

WL is a widely used time-domain feature for EMG-based solutions since it efficiently combines amplitude and frequency information~\cite{HandGestureRecognition, 
ComprehensiveStudyOnEMG, BiomedicalSignalProcessingAndControl, Review_IRBM_EMG, ReviewOfEMGSignalClassification}.

\subsubsection{Peak Frequency - PKF} 
\label{sssec:peak_frequency}

Peak Frequency is the frequency where the maximum power occurs in the signal's power spectrum.
The PKF is obtained by finding the frequency at which the power spectral density reaches its maximum value:

\begin{equation}
  PKF = \arg\max_f P(f)
  \label{eq:pkfEQ}
\end{equation}

Where $P(f)$ is the power spectral density of the signal at frequency $f$.
This feature provides valuable information about the dominant frequency of muscle activity, which is often related to the type and intensity of the contraction. Changes in PKF 
can indicate muscle fatigue, as the dominant frequency tends to shift towards lower values when fatigue increases. Therefore, PKF is commonly used in EMG analysis to assess 
muscle performance and fatigue levels~\cite{HandGestureRecognition, Review_IRBM_EMG}.

\subsubsection{Median Frequency - MDF} 
\label{sssec:median_frequency}

The Median Frequency corresponds to the frequency where the power spectrum is equally divided in two. This can be translated to the following equation:

\begin{equation}
  MDF = \frac{1}{2}\sum_{i = 1}^{N} P_i
  \label{eq:mdfEQ}
\end{equation}

Where $P_i$ is the power spectral density of the signal at frequency $i$, and $N$ is the maximum frequency component considered~\cite{HandGestureRecognition, Review_IRBM_EMG, 
ReviewOfEMGSignalClassification, ComprehensiveStudyOnEMG}.


\subsubsection{Mean Frequency - MNF} 
\label{sssec:mean_frequency}

Mean Frequency represents the average frequency of the power spectrum at each frequency component, weighted by the power spectral density.

\begin{equation}
  MNF = \frac{\sum_{i = 1}^{N} f_i P(f_i)}{\sum_{i = 1}^{N} P(f_i)}
  \label{eq:mnfEQ}
\end{equation}

Where $f_i$ represents the $i^{th}$ frequency component, $P(f_i)$ is the power spectral density at that frequency, and $N$ is the total number of frequency bins considered.

MNF provides a good insight into the distribution of power across the frequency spectrum and is often used to assess muscle fatigue, since it tends to decrease as fatigue increases. 
Unlike MDF~(\autoref{sssec:mean_frequency}), which divides the spectrum into two equal halves, MNF reflects the overall spectral balance of the signal~\cite{HandGestureRecognition, 
ComprehensiveStudyOnEMG, Review_IRBM_EMG, ReviewOfEMGSignalClassification, Hybrid, EvaluationOfFeatureExtractionAndRecognition}.


%\subsubsection{Wavelet Packet Energy - WPE} % HandGestureRecognition
%\label{sssec:wavelet_packet_energy}

\subsubsection{Integrated EMG - IEMG}
\label{sssec:integrated_emg}

Integrated~\gls{EMG} is the sum of the rectified~\gls{EMG} signal over a time interval. This feature quantifies the muscle activation level providing information about the effort 
and muscle intensity during a movement.

Mathematically, the integrated EMG is defined by:

\begin{equation}
  IEMG = \sum_{i = 1}^{N} \left\lvert x_i \right\rvert 
  \label{eq:iemgEQ}
\end{equation}

Where $x_i$ is the~\gls{EMG} signal amplitude at each sample, and $N$ is the total number of samples in the analyzed time period.
IEMG is particularly useful for assessing the intensity of muscular activity and comparing activation levels across different tasks or conditions~\cite{ComprehensiveStudyOnEMG, 
BiomedicalSignalProcessingAndControl, Review_IRBM_EMG}.

\subsubsection{Simple Square Integral - SSI} 
\label{sssec:simple_square_integral}

The Simple Square Integral is a feature that expresses the energy of the~\gls{EMG} signal. 
The equation below expresses the computation of SSI:

\begin{equation}
  SSI = \sum_{i = 1}^{N} \left\lvert x_i^2 \right\rvert 
  \label{eq:ssiEQ}
\end{equation}

Where $x_i$ represents the~\gls{EMG} signal amplitude at the $i^th$ sample and N is the total number of samples within the analysis window\cite{ComprehensiveStudyOnEMG, 
BiomedicalSignalProcessingAndControl, Review_IRBM_EMG}.

\subsubsection{Zero Crossing - ZC} 
\label{sssec:zero_crossing}

Zero Crossings is a feature that measures the amount of times that the~\gls{EMG} signal changes sign as expressed by the following expression:

\begin{equation}
  ZC = \begin{cases}
    \sum_{n = 2}^{N} sgn(x_i \cdot x_{i-1})
    (x_i - x_{i-1})~\geq \epsilon 
  \end{cases}
  \label{eq:zcEQ}
\end{equation}
  Where sgn($\cdot$) is the sign function

\begin{equation}
  sgn(x) = \begin{cases}
    1 & \text{if }  x > \epsilon \\
    0 & \text{otherwise}
  \end{cases}
  \label{eq:zcEQ2}
\end{equation}

Where $x_i$ is the~\gls{EMG} signal amplitude at each sample, $N$ is the total number of samples and $\epsilon$ is a small positive value used to avoid counting zero crossings 
caused by minor fluctuations or noise\cite{Review_IRBM_EMG}.

\subsubsection{Slope Sign Changes - SSC}
\label{sssec:slope_sign_changes}

The Slope Sign Change feature measures the number of times the slope of the EMG signal changes sign within a given time window. It reflects the signal's frequency and 
complexity, providing information about the motor unit firing pattern and the muscle's dynamic behaviour. A higher number of slope sign changes generally indicates a more complex or higher-frequency muscle activity.

As with the ZC~(\autoref{sssec:zero_crossing}) feature, a small threshold is often used to minimize the effect of noise on the calculation. 

Mathematically, the SSC can be expressed as:

\begin{equation}
  SSC = \sum_{i = 2}^{N - 1} f[(x_i - x_{i-1}) \cdot (x_i - x_{i+1})]
  \label{eq:sscEQ}
\end{equation}

Where

\begin{equation}
  f(x) = \begin{cases}
    1 & \text{if } x~\geq \epsilon \\
    0 & \text{otherwise}
  \end{cases}
  \label{eq:sscEQ2}
\end{equation}

Where, $x_i$ represents the EMG signal amplitude at sample $i$, and the $\epsilon$ is a small positive value used to reduce sensitivity to minor fluctuations or 
electrical noise~\cite{Review_IRBM_EMG}.

\subsubsection{Willison Amplitude - WAMP}
\label{sssec:wamp}

Willison Amplitude is a time-domain feature that assumes the number of times the absolute difference between two consecutive signal samples exceeds a predefined 
threshold. It is particularly useful for capturing changes in muscle activation intensity while remaining computationally efficient\cite{Review_IRBM_EMG, 
EvaluationOfFeatureExtractionAndClassification, EvaluationOfFeatureExtractionAndRecognition, ComparativeStudyOfSEMGFeatureEvaluationMethods, ResearchOnLowerLimbMotionRecognition}.

Mathematically, it is defined by:

\begin{equation}
  WAMP = \sum_{n = 1}^{N - 1} f(\left\lvert x_i - x_{i+1}\right\rvert )
  \label{eq:wampEQ}
\end{equation}

Where

\begin{equation}
  f(x) = \begin{cases}
    1 & \text{if } x~\geq \epsilon \\
    0 & \text{otherwise}
  \end{cases}
  \label{eq:wampEQ2}
\end{equation}

\section{Related Work}
\label{sec:related_work}

For sports like basketball, the available motion recognition research is mainly through video recording or smart wearable devices such as 
wristbands. There are many studies on the topic of measuring plantar pressure through an insole. Most of them mostly study the gait cycle 
while running or walking with similar objectives.

\subsection{Muscle Cramps and Pre-Cramp Anomalies}
\label{ssec:cramps}

Cramping occurring in sports or conditioning is clinically called an~\gls{EAMC}. This is a sudden, involuntary, 
and painful contraction of skeletal muscle during or immediately following physical activity.~\glspl{EAMC} severely 
disrupt athletic performance and can precipitate further musculoskeletal injury. While historically attributed to 
systemic factors such as dehydration or electrolyte imbalance, contemporary sports medicine increasingly supports 
the "altered neuromuscular control" hypothesis. This theory posits that severe, localized muscle fatigue disrupts 
the neural mechanisms regulating muscle spindle and Golgi tendon organ activity, leading to sustained, involuntary 
motor unit firing. Whenever an EAMC happens, the EMG signal demonstrates an extreme and abnormal behaviour. The 
focal point of the cramp is characterized by massive localized spikes in~\gls{sEMG} amplitude~\cite{PossibleMechanisms}. During the 
cramping, the~\gls{RMS} and~\gls{MAV} features record very high values when compared to standard movements. 

Current research emphasizes that acute muscular failure is rarely an instantaneous event but rather precedented by 
a quantifiable period of continuous, localized muscular fatigue and deteriorating neuromuscular efficiency. The 
best and most common approaches for detecting this fatigue utilize~\gls{sEMG} to track its temporal evolution. 
Physiological changes, such as the accumulation of metabolic byproducts slow down muscle fibre conduction velocity, 
fundamentally altering the~\gls{sEMG} signal~\cite{TheUseOfSEMG}. Time-domain features are highly effective markers for this 
phenomenon. As muscles fatigue, amplitude-based features such as RMS and MAV typically exhibit an upward trend as 
the central nervous system increases motor unit recruitment and synchronization to maintain mechanical force. On 
the other hand, features that act as proxies for frequency, such as~\gls{ZC} and~\gls{SSC} decrease as action 
potentials slow down. Research has successfully utilized these specific time-domain features to construct objective 
fatigue indices, proving their efficiency in capturing impending muscle failure~\cite{ProposedFatigueIndex}. 

Recent literature shows that 
these involuntary contractions aren't spontaneous but rather have some distinct pre-cramp anomalies. For instance, 
in pilot studies observing involuntary nocturnal leg cramps, researchers found that changes in muscle tension 
patterns couldreliably predict the start of the clinical cramp. Specifically, sEMG amplitude was observed spiking 
significantly, in some cases up to 40 seconds before the painful contraction physically manifested or triggered 
any movement artifacts~\cite{ChangesInMuscleTension}.

\subsection{Sensor Fusion}
\label{ssec:sensor_fusion}


Despite its proven efficacy, relying on isolated sEMG data presents critical limitations in dynamic athletic 
environments. A high sEMG amplitude indicates elevated electrical demand, which could signify impending fatigue-induced 
failure, or it could simply reflect an athlete's intentional increase in mechanical output, such as executing a 
jump or rapidly accelerating. To resolve this ambiguity, modern systems utilize sensor fusion, coupling~\gls{sEMG} 
with the~\glspl{IMU} discussed previously. By correlating the electrical demand of the muscle (via~\gls{sEMG}) with 
the actual mechanical output or kinematic stability (via~\gls{IMU}), researchers can accurately map true neuromuscular
efficiency. Recent studies confirm that multimodal fatigue detection systems using both kinematic and neuromuscular 
signals provide highly robust continuous fatigue classification, achieving accuracies above 95\% even in the presence 
of complex, continuous movements~\cite{DataDrivenApproachForUpper, MultimodalFatigueDetectionSystem}. 

Edge computing and machine learning algorithms, such as~\glspl{RF} and Decision Trees, represent the cutting edge 
of real-time monitoring~\cite{MachineLearningBasedClassification}. By identifying the precise moment neuromuscular efficiency begins to decline, these 
unobtrusive systems provide the foundational architecture necessary to predict and prevent specific, acute fatigue-related 
conditions. 

Qingsong Ai et al. developed a lower limb motion recognition system that successfully classified five distinct 
movements by fusing~\gls{EMG} and accelerometer signals. To handle the kinematic data, the researchers applied a 
curve fitting method based on median filtering to remove noise, and subsequently utilized~\gls{DTW} to extract 
sequential distance features. Simultaneously, the~\gls{sEMG} signals were processed using a power spectral correlation 
coefficient method for accurate movement onset detection, followed by the extraction of both time-domain features 
and wavelet transform coefficients. By inputting these fused multimodal features into a~\gls{SVM} classifier 
optimized via grid search, the system achieved an average recognition accuracy of 98.1\%. Ultimately, this research 
proved that combining neuromuscular data with kinematic accelerometer features significantly outperforms classification 
models that rely on either signal modality in isolation~\cite{ResearchOnLowerLimbMotionRecognition}. 

In another study focusing on the integration of upper and lower limb kinematics, Bárbara Cardoso explored a multimodal 
wearable system combining smart insoles with a Moto 360 smartwatch to optimize sports performance, specifically in 
running. The methodology relied on extracting a robust set of 54 time- and frequency-domain features from the smartwatch's 
triaxial accelerometer. To ensure computational efficiency for realtime edge application,~\gls{RFE} was employed to 
isolate the top 10 most discriminative features, which predominantly included the~\gls{RMS}, StandardDeviation, and 
\gls{ZC} Rate of the x-axis acceleration. These features were subsequently fed into a Multilayer Perceptron~\gls{NN}, 
achieving an impressive 95\% overall accuracy across five distinct physical activities, with running specifically 
classified at a 90\% accuracy and 97\% precision. Beyond activity recognition, the system utilized the smartwatch's 
photoplethysmography sensor to estimate Energy Expenditure and successfully demonstrated that arm acceleration 
amplitude varies in indirect proportion to foot ground contact time, establishing a clear biomechanical link 
between wrist-worn sensor data and lower-limb running dynamics~\cite{porto}.

\subsection{Insole and Pressure Sensors}
\label{ssec:insole_preassure}

Baitong Wang et al. developed an insole with 8 pressure sensors and motion track sensors and with this system managed to clearly identify 
gait with only the pressure sensors~\cite{freewalker}.

Elizabeth S. Chumanov et al. tracked the global position of the pressure on the soles of the feet throughout steps in walking and running. 
This data can be used for improvement in foot-floor contact models used to simulate human locomotion and in least squares inverse dynamics
~\cite{posDuringWalking}. 

Thomas Holleczek et al. design a pair of smart socks with three pressure sensors each for the use case of snowboard. Even tho it isn't a 
sport as dynamic and unpredictable as basketball, it still is very relevant for this project. They worked on activity recognition focusing 
on recognizing if the user was riding backside or riding frontside (with his back to the bottom of the hill or the top respectively). It is 
pointed that the recognizability of three basic standing postures was explored and that they were distinguished with perfect accuracy 
mentioning that it can be interesting for gait analysis~\cite{snowboard}. 

Qing F. Zhou et al. utilized insoles with a motion sensor with the purpose of getting a high accuracy for recognizing three basketball 
motions: Dribbling, jumping and turning around. For this data processing the K-means algorithm was utilized~\cite{basketballMotions}. Two years later they 
continued their work and managed to identify five new different foot work movements. In this version a multiple sensor system was proposed 
and three different classification methods were discussed:~\gls{KNN},~\gls{SVM} and~\gls{CNN}~\cite{basketballFootwork}. 

Catarina Amaro et al. studied the differences in five different movements in basketball between two different moments in a season. The 
study presented no significant changes in mean peak pressure between the two different times but it did specify that experienced players 
had lower plantar pressure values in their movements and lower symmetry index. This means that there is room for improvement based on the 
values read for the less experienced players. It was also mentioned that wear of the shoes may have influenced the slight changes observed 
in the study.Each player wore their own shoes and this was pointed as a variable for the plantar pressure measured~\cite{coimbra}.


\subsection(Summary)
\label{ssec:rw_summary}

\begin{table*}[ht]
\centering
\caption{Summary of Related Work in Activity Recognition, Sensor Fusion, and Fatigue Monitoring.}
\label{tab:related_work}
\renewcommand{\arraystretch}{1.2}
\footnotesize
\begin{tabularx}{\textwidth}{|p{2cm}|p{2cm}|p{2cm}|p{2cm}|X|}
\hline
\textbf{Reference} &
\textbf{Sensors / Modality} &
\textbf{Application / Focus} &
\textbf{Methods / Algorithms} &
\textbf{Key Results} \\
\hline

Ai et al.~\cite{ResearchOnLowerLimbMotionRecognition} &
sEMG + Accelerometer &
Lower limb motion recognition (5 movements) &
\gls{DTW},~\gls{SVM} (grid search optimization) &
Achieved 98.1\% accuracy. Proved sensor fusion significantly outperforms isolated sensors. \\
\hline

Cardoso~\cite{porto} &
Smartwatch (\gls{IMU}, Photoplethysmography) + Smart Insoles &
Running / Sports performance optimization &
\gls{RFE}, Multilayer Perceptron~\gls{NN} &
95\% overall activity recognition accuracy; established biomechanical link between arm acceleration and ground contact. \\
\hline

Wang et al.~\cite{freewalker} &
8 Pressure Sensors + Motion Trackers &
Gait identification &
-- &
Successfully identified gait patterns relying solely on plantar pressure data. \\
\hline

Chumanov et al.~\cite{posDuringWalking} &
Plantar Pressure Sensors &
Walking and running gait analysis &
Least squares inverse dynamics &
Tracked global pressure positions to improve simulated human locomotion and foot-floor contact models. \\
\hline

Holleczek et al.~\cite{snowboard} &
Smart Socks (3 Pressure Sensors) &
Snowboard activity recognition &
-- &
Achieved perfect accuracy in distinguishing three basic standing postures; highlighted potential for gait analysis. \\
\hline

Zhou et al.~\cite{basketballMotions} &
Smart Insoles + Motion Sensor &
Basketball motion recognition (3 movements) &
K-means Clustering &
High accuracy achieved for recognizing dribbling, jumping, and turning around. \\
\hline

Zhou et al.~\cite{basketballFootwork} &
Multiple Sensors (Insoles) &
Basketball footwork recognition (5 movements) &
\gls{KNN},~\gls{SVM},~\gls{CNN} &
Successfully classified complex footwork by comparing traditional~\gls{ML} against deep learning. \\
\hline

Amaro et al.~\cite{coimbra} &
Plantar Pressure Sensors &
Plantar pressure in basketball (5 movements) &
Symmetry Index Analysis &
Found that more experienced players exhibit lower plantar pressure values and a lower symmetry index. \\
\hline

Qassim et al.~\cite{ProposedFatigueIndex} &
Isolated sEMG &
Muscle fatigue detection &
Time-domain analysis (\gls{RMS},~\gls{MAV},~\gls{ZC},~\gls{SSC}) &
Demonstrated that amplitude features increase and frequency features decrease, efficiently capturing impending muscle failure. \\
\hline

Chaney et al.~\cite{ChangesInMuscleTension} &
Isolated~\gls{sEMG} &
\gls{EAMC} (cramp) prediction &
Muscle tension pattern observation &
Identified that~\gls{sEMG} amplitude spikes up to 40 seconds before clinical cramps, revealing a window for proactive prevention. \\
\hline

\end{tabularx}
\end{table*}




% \printbibliography[heading=subbibliography, segment=\therefsegment, title={\bibname\ for chapter~\thechapter}]

